\section{Conclusion}

An evaluation cascade applied to 225 model configurations across three genotoxicity datasets 
shows that, for the Ames mutagenicity benchmark studied here, data source 
heterogeneity---the mixing of drug-like and industrial compounds with sharply different 
prevalence and provenance---drives a substantially larger performance gap ($\Delta = 0.390$, 
scaffold CV $\to$ LDO) than the choice between random and scaffold splitting 
($\Delta = 0.046$). A naive source classifier using no structural information reaches MCC = 
0.608 on the same benchmark, suggesting that mixed-source Ames performance partly reflects 
learned source membership rather than mutagenicity. Threshold-tuning analysis attributes roughly one-third (30.9--35.3\,\%) of the LDO gap to 
prior shift, leaving a residual that no recalibration can fix. 
Algorithm differences are small in-domain (MCC range = 0.063) but become pronounced under 
source shift, with RF more robust than gradient boosting in the harder direction; this 
cross-source advantage is feature-conditional and warrants further confirmation. In-vivo 
micronucleus models show ranking signal (PR-AUC 8.5-fold above baseline) but classification 
metrics whose CIs do not exclude zero on the present test set. These findings motivate an 
evaluation cascade with source-awareness as a useful complement to existing practices for 
mutagenicity QSAR.
